\chapter{نتیجه‌گیری}
\newpage
‎\renewcommand{\baselinestretch}{2}\normalsize‎
\label{Conclusion}

همانطور که در این رساله بیان کردیم، رویکرد یادگیری برای شبکه‌های تطبیقی توزیع‌شده ناهمگن طراحی و ارائه نشده است. سایر روش‌های یادگیری توزیع‌شده مطرح شده یا به صورت متمرکز پیاده‌سازی شده‌اند و یا شبکه را به صورت مجموعه‌ای از گره‌های همگن فرض کرده‌اند. هر کدام از این روش‌ها نقاط ضعفی دارند که در فصل‌های قبل بیان کرده‌ایم. در این فصل به بیان نتیجه‌گیری کل کار پرداخته می‌شود و در پایان نیر به بیان پیشنهادهایی در راستای ادامه این کار خواهیم پرداخت. 

\section{نتیجه‌گیری}
امروزه روش‌های موازی‌سازی و مدل‌های محاسباتی به دلیل افزایش حجم محاسبات و داده‌های پردازشی در بسیاری از زمینه‌ها و کاربردها توجه زیادی را به خود جلب نموده‌اند. همچنین با پیشرفت روز افزون تکنولوژی و قدرت‌های پردازشی متفاوت سیستم‌ها، فرایند یادگیری در شبکه‌‌ای از گره‌های پردازشی با محدودیت‌ها و چالش‌های مهمی مواجه خواهد بود. در یک شبکه تطبیقی توزیع‌شده، اگر عملیات محاسباتی برای گره‌ها سنگین باشد و یا داد‌های ورودی ذاتاً توزیع‌شده و سنگین باشند، مهمترین چالش در این نوع یادگیری، موازی‌سازی فرآیند آموزش و ایجاد یک مدل منسجم در کل شبکه است.  از آنجایی که موازی‌سازی بهتر است تحت یک ساختار و مدل مشخص و کارایی انجام شود، ارائه یک مدل محاسباتی موازی منطبق با ویژگی‌های شبکه‌های تطبیقی توزیع‌شده با قابلیت اجرای کارا در محیط‌های ناهمگن بسیار چالش برانگیز بوده اما دستیابی به آن هدفی بسیار ارزشمند است. 

مدل محاسباتی موازی \lr{BSP} مدلی مناسب برای شبکه‌های تطبیقی توزیع‌شده می‌باشد اما کلیه گره‌های شرکت کننده در فرایند یادگیری این شبکه در سطح یکسانی قرار دارند. از آنجایی که گره‌های شبکه ممکن است از لحاظ قدرت محاسباتی و پردازشی متفاوت باشند، مدل محاسباتی \lr{BSP} برای این حالت مناسب و کارا نمی‌باشد. از این رو در این رساله، به طراحی و ارائه مدل محاسباتی توسعه‌یافته‌ای به نام \lr{BSP} فدرالی پرداخته شد. طراحی چنین مدلی با هدف کاهش زمان انتظار گره‌ها در طول فرایند یادگیری و کاهش میزان ارتباطات با حفظ سرعت همگرایی و دقت یادگیری انجام شد. 

رویکرد یادگیری توزیع‌شده مبتنی بر مدل محاسباتی \lr{BSP} فدرالی ارائه شده از دو مرحله اصلی برون خط و برخط تشکیل شده است. در مرحله اول گره‌های شبکه در منطقه‌هایی مشابه خوشه‌بندی می‌شوند. در ساختار مدل محاسباتی توسعه یافته ارتباطات به منظور تبادل اطلاعات و همچنین همگام‌سازی‌های مانعی به منظور همگام‌شدن تخمین‌های گره‌ها در انتهای هر ابرگام بر اساس منطقه‌ها انجام می‌شود. در مرحله دوم فرایند یادگیری توزیع‌شده توسط گره‌های منطقه‌بندی‌شده به دو صورت منطقه‌ای و سرتاسری انجام می‌شود. از آنجایی که گره‌ها با قدرت پردازشی مشابه در یک منطقه قرار می‌گیرند، گره‌های کند دیگر مانع همگام‌سازی دیرهنگام گره‌های سریع‌تر نمی‌شوند. در ابرگام‌های منطقه‌ای، هر منطقه مستقل از سایر منطقه‌ها فرایند یادگیری را به صورت تطبیقی توزیع‌شده اجرا می‌کند. در ابرگام‌های سرتاسری، هر گره فرایند یادگیری را به صورت عادی و در سطح کل شبکه انجام می‌دهد. همین امر منجر به کاهش زمان انتظار گره‌ها در زمان اجرا فرایند تخمین پارامترها شد و همچنین منجر به کاهش تعداد ارتباطات و تبادل اطلاعات بین گره‌ها گردید. 

\section{پیشنهادات آتی}
اگر چه نتایج حاصل از این پژوهش رضایت‌بخش و قابل قبول می‌باشد، با این حال جنبه‌های جذاب و جالب توجهی پیش‌روی این زمینه پژوهشی وجود دارد. می‌توان مباحث امنیت داده‌ها را در این کار دخیل کرد و کار تحقیقاتی جدید را تعریف نمود. همچنین می‌توان رویکرد پیشنهادی را برای شبکه‌های عصبی عمیق استفاده کرد و آن را در این زمینه توسعه داد و راهکار موثری عرضه کرد. از سوی دیگر می‌توان مدل‌های محاسباتی دیگر را نیز به صورت فدرالی توسعه داد و نسخه‌های جدیدی از آنها برای کار بر روی داده‌های غیرمستقل با توزیع‌های متفاوت عرضه کرد.

\section{مقالات مستخرج از رساله}
با انجام این رساله، دو مقاله چاپ شده و یک مقاله ارسال شده از آن استخراج شده است. مقاله اول مربوط به الگوریتم یادگیری توزیع‌شده پیشنهادی می‌باشد که نتیجه آن چاپ یک مقاله در مجله \lr{Journal Signal Processing} با عنوان \lr{‌‌‌Convergence Behavior of Diffusion Stochastic Gradient Descent Algorithm} در سال 2021 شده است. مقاله دوم مربوط به مقاوم‌سازی الگوریتم یادگیری پیشنهادی و مساله پیش‌بینی سیگنال‌های بازار بورس می‌باشد که نتیجه آن چاپ یک مقاله در مجله \lr{Journal of Parallel and Distributed Computing} با عنوان \lr{A Distributed Learning based on Robust Diffusion SGD over Adaptive Networks with Noisy Output Data} در سال 2024 شده است. مقاله سوم مربوط به مدل محاسباتی پیشنهادی می‌باشد که به مجله \lr{Journal of IEEE Transactions on Parallel and Distributed Systems} ارسال شده است و تحت داوری است. 
%\renewcommand{\arraystretch}{2}
%\setlength{\tabcolsep}{25t}
\begin{table}[h]
	\centering
	\caption{\small{لیست مقالات مستخرج از رساله}}
	\footnotesize\rm
%	\resizebox{\textwidth}{!}{
	\begin{tabular}{|c|p{4cm}|p{4cm}|p{4cm}|c|}
		\hline
		ردیف&عنوان&نویسندگان&محل انتشار&تاریخ\\
		\hline
	1&\lr{‌‌‌Convergence Behavior of Diffusion Stochastic Gradient Descent Algorithm}&\lr{Fatemeh Barani, Abdorreza Savadi, Hadi Sadoghi Yazdi}&\lr{Journal Signal Processing}&2021\\
	\hline
	2&\lr{A Distributed Learning based on Robust Diffusion SGD over Adaptive Networks with Noisy Output Data}&\lr{Fatemeh Barani, Abdorreza Savadi, Hadi Sadoghi Yazdi}&\lr{Journal of Parallel and Distributed Computing}&2024\\
	\hline
	3&\lr{Federated Bulk Synchronous Parallel Model for Distributed Adaptive Networks}&\lr{Fatemeh Barani, Abdorreza Savadi, Hadi Sadoghi Yazdi}&\lr{Journal of IEEE Transactions on Parallel and Distributed Systems}&\lr{Submit}\\
		\hline
	\end{tabular}
%}
\end{table}

\section{خلاصه}
در این فصل تحلیل نتایج و نقاط قوت و ضعف رویکرد پیشنهادی بیان شده است و با توجه به نقاط ضعفی که در رویکرد پیشنهادی وجود دارد، پیشنهادهایی نیز برای مطالعات آینده بیان شد. همچنین مقالات مستخرج شده از این رساله آورده شده‌اند. 